% !TeX root = ../../Book2.tex
\chapter{Poisson 求和公式与 \texorpdfstring{\(\theta\)}{\AnalysisPdfTheta} 函数预备}
\label{b2:ch:06}
\BookChapterNumber{b2:ch:06}
% Hunter2001Applied Section11.11印刷319页实际图像核读：11.45与Example11.46。
% 本章固定单位周期及2π频率约定，完整展开周期化换序、尺度、格与对偶格。
\begin{chapterabstract}
Poisson 求和公式通过周期化函数的 Fourier 系数，把格点求和问题化为对偶格上的求和。本章从光滑周期化出发证明公式，先讨论缩放求和与数值求积，再通过 Gauss 函数得到尺度倒置与圆周热核的两种表达。带限采样与频率混叠单列为选读专题，PDE 路线可以跳过。实参数 \(\theta\) 级数随后把这种对称连接到格点计数和整数平方和。高维选读部分讨论格基、对偶格与基本区域体积，为第三卷研究 \(\theta\) 函数及解析数论作准备。
\end{chapterabstract}
\begin{learninggoals}
\begin{itemize}
\item\label{b2:item:06-goal-periodic}从周期化函数计算 Fourier 系数，逐项核对求和、积分和微分的交换条件。
\item\label{b2:item:06-goal-scale}写出带平移、调制和尺度的 Poisson 公式，解释格点间距与频率间距的倒数关系，分别控制数值求积的频率误差与空间尾部；选读\subsecref{b2:subsec:060105}进一步研究带限恢复和混叠。
\item\label{b2:item:06-goal-theta}由 Gauss 自对偶性推导 \(\theta\) 变换，使用指数尾在大小参数之间转换估计。
\item\label{b2:item:06-goal-count}区分平滑格点计数与球内硬计数，理解平方和生成函数及 Mellin 积分的作用。
\item\label{b2:item:06-goal-lattice}在选读部分计算对偶格与基本区域体积，保留高维公式中的行列式因子。
\end{itemize}
\end{learninggoals}

怎样用等间距的数值 \(f(nh)\) 估计整条实线上的积分？减小 \(h\) 会改善近似，但它也改变频率中被读出的点。对 Gauss 函数，这个误差能精确改写成另一个快速收敛的级数；对带限信号，离散点值甚至能够完全恢复函数。本章以周期化把这些问题放进同一个计算：空间中的格点求和，对应频率中的倒数格点求和。

\ProseParagraphBreak

本章同时使用\chapref{b2:ch:02}的周期系数唯一性与\chapref{b2:ch:05}的 Schwartz 变换不变性。Gauss 变换沿用\chapref{b2:ch:03}已经固定的常数；平方能量观点则帮助解释热核的频率表达。求和公式的证明同时使用离散展开与连续变换。

\ProseParagraphBreak

若以 PDE 为目标，先读\secref{b2:sec:0601}的周期化、Poisson 求和及\subsecref{b2:subsec:060104}的缩放公式，再读\secref{b2:sec:0602}的 Gauss 函数与热核。对离散测量如何恢复连续信号感兴趣的读者，可继续读\subsecref{b2:subsec:060105}的带限采样和混叠；Hermite 自对偶补例则展示更多变换下保持不变的函数。\secref{b2:sec:0603}的平方和计数与 Mellin 算术联系留给需要连接第三卷的读者。\secref{b2:sec:0604}将整数格推广到一般格，说明格的形状与对偶格怎样进入求和公式。John K. Hunter 与 Bruno Nachtergaele 的《Applied Analysis》\cite[Section 11.11, p. 319]{Hunter2001Applied}以周期化证明求和公式，并给出 theta 级数的例子；Fourier 分析在数论中的应用还见于 Elias M. Stein 与 Rami Shakarchi 的《Fourier Analysis: An Introduction》\cite{Stein2003Fourier}。

% !TeX root = ../../Book2.tex
\section{二阶线性方程}
\label{b2:sec:0801}

两个二阶方程可能只差一个符号，解的行为却截然不同。分类首先抓住最高阶导数的组合，把它看成每一点上的二次型。本节使用 \(z=(z_1,\ldots,z_n)\) 表示全部自变量；只有在讨论演化问题时才拆成 \(z=(t,x)\)。因此 \(n\) 不总等于空间维数。

\subsection{主部}
\label{b2:subsec:080101}

在开集 \(\Omega\subseteq\mathbb R^n\) 上考虑实系数算子
\begin{equation}\label{b2:eq:08-linear-operator}
 Lu=\sum_{i,j=1}^n a^{ij}(z)\partial_i\partial_ju
 +\sum_{i=1}^n b^i(z)\partial_i u+c(z)u.
\end{equation}
系数至少连续，经典解至少属于 \(C^2\)。一阶和零阶项称为低阶项；主部已在\secref{b2:sec:0104}引入，此处专指含二阶导数的和式，记为 \(L_2u\)。由于混合偏导相等，把 \(a^{ij}\) 换成 \((a^{ij}+a^{ji})/2\) 不改变算子。以下始终取实对称矩阵 \(A=(a^{ij})\)。

\ProseParagraphBreak

二维中习惯写为
\[
 L_2u=a u_{xx}+2b u_{xy}+c u_{yy},\quad
 A=\begin{pmatrix}a&b\\b&c\end{pmatrix}.
\]
混合项前面的 \(2\) 来自矩阵求和中 \((1,2)\) 和 \((2,1)\) 两项。若方程写成 \(u_{xx}+6u_{xy}+u_{yy}=0\)，则矩阵的非对角元为 \(3\)，不是 \(6\)。本节中的 \(a,b,c\) 是二维主部的临时系数，不与\eqref{b2:eq:08-linear-operator}中的向量低阶系数混用。

\ProseParagraphBreak

散度形式也有相同的最高阶矩阵。在 \(a^{ij}\in C^1\) 时，
\[
 \sum_{i,j}\partial_i(a^{ij}\partial_ju)
 =\sum_{i,j}a^{ij}\partial_i\partial_ju
 +\sum_{i,j}(\partial_i a^{ij})\partial_ju.
\]
它与非散度形式在经典分类上相同，但分部积分时低阶项不同；在粗糙系数或弱解问题中，两种写法更不能只凭主部相同便随意交换。

\subsection{主符号}
\label{b2:subsec:080102}

\begin{definition}\label{b2:def:08-principal-symbol}
设 \(n\geq1\)，\(\Omega\subset\mathbb R^n\) 为开集，实系数算子 \(L=\sum_{i,j}a^{ij}(z)\partial_{ij}+\sum_i b^i(z)\partial_i+c(z)\) 的系数均连续，且 \(A=(a^{ij})\) 为实对称矩阵场。对 \(z\in\Omega\)，二次齐次多项式
\[
 p(z,\xi)\coloneqq\xi^{\mathsf T}A(z)\xi
 =\sum_{i,j=1}^n a^{ij}(z)\xi_i\xi_j
 \quad(\xi\in\mathbb R^n)
\]
称为 \(L\) 的\term[zhufuhao]{主符号}[principal symbol]。本章不把 Fourier 微分规则中的 \(-4\pi^2\) 因子纳入这个几何记号。
\end{definition}
\symidx{\(p(z,\xi)\)：二阶算子的几何主符号}

后一约定与第一编的 \(2\pi\) 归一化相容：对常系数主部，
\[
 L_2e^{2\pi\mathrm i z\cdot\xi}
 =-4\pi^2p(\xi)\cdot e^{2\pi\mathrm i z\cdot\xi}.
\]
因此 \(p\) 记录不同频率方向的二阶作用，整个 Fourier 乘子则是 \(-4\pi^2p\)。不同文献可能把这个负号纳入符号定义；判断椭圆性和特征集合时，必须先核对其约定。

\ProseParagraphBreak

更一般地，对光滑实相位 \(\phi\) 与大参数 \(\lambda\)，直接两次求导给出
\[
 L\bigl(e^{2\pi\mathrm i\lambda\phi}\bigr)
 =e^{2\pi\mathrm i\lambda\phi}
 \left(-4\pi^2\lambda^2p(z,\nabla\phi)
 +2\pi\mathrm i\lambda\left(\sum_{i,j}a^{ij}\partial_{ij}\phi
 +\sum_i b^i\partial_i\phi\right)+c\right).
\]
在固定紧集上，最高的 \(\lambda^2\) 次项由主符号决定。若该系数为零，通常占主导的二阶作用恰好抵消；特征几何正由这类方向产生。这个计算只是代数恒等式，没有断言任意高频函数就是某个方程的近似解。

\subsection{特征余向量}
\label{b2:subsec:080103}

\begin{definition}\label{b2:def:08-characteristic-covector}
设 \(n\geq1\)，\(\Omega\subset\mathbb R^n\) 为开集，实系数算子 \(L=\sum_{i,j}a^{ij}(z)\partial_{ij}+\sum_i b^i(z)\partial_i+c(z)\) 的系数均连续，且 \(A=(a^{ij})\) 为实对称矩阵场。令 \(p(z,\nu)=\nu^{\mathsf T}A(z)\nu\)。在一点 \(z\in\Omega\)，若非零实余向量 \(\nu\in(\mathbb R^n)^*\) 满足 \(p(z,\nu)=0\)，则称它为该点的\term[tezheng-yuxiangliang]{特征余向量}[characteristic covector]。若实值函数 \(\phi\in C^2(V)\) 定义在开集 \(V\subseteq\Omega\) 上，在超曲面 \(\Sigma=\{z\in V:\phi(z)=0\}\) 上处处有 \(\nabla\phi\ne0\)，且 \(p(z,\mathrm d\phi(z))=0\)，则称 \(\Sigma\) 为\term[tezheng-chaqumian]{特征超曲面}[characteristic hypersurface]。在 \(\Sigma\) 上满足 \(p(z,\mathrm d\phi(z))\ne0\) 的点称为该曲面的非特征点。
\end{definition}

这里的 \(\nu\) 表示余向量，也就是对位移作线性求值的法向数据；借欧氏内积可用列向量表示，但不能把它直接当作曲面的切向。例如，\(\mathrm d\phi(z)\) 对切向量 \(w\) 的作用是 \(\mathrm d\phi(z)[w]=\nabla\phi(z)\cdot w\)，在切平面上等于零。下文计算主符号时，以梯度的列向量表示同一个微分。非零倍数 \(\alpha\nu\) 仍为特征余向量，所以也可以讨论它确定的、忽略非零实倍数的余法向方向；二维特征曲线的切方向则是另一个对象。

\ProseParagraphBreak

换一个定义函数 \(\widetilde\phi=\alpha\phi\)，其中 \(\alpha\) 是曲面附近不为零的 \(C^2\) 实函数，在曲面上有 \(\mathrm d\widetilde\phi=\alpha\,\mathrm d\phi\)。主符号齐次，零值条件不变。一般地，同一光滑超曲面的两个正则定义函数在曲面上的非零微分都有同一个核，即切平面，因而成非零比例；所以特征条件不依赖定义函数的选择。一个曲面也可能只有部分点为特征点，不能从一个点的条件推出整张曲面均为特征。

\ProseParagraphBreak

在二维中，若曲线切向为 \((\mathrm dx,\mathrm dy)\)，则其一个法向与 \((-\mathrm dy,\mathrm dx)\) 平行。特征条件因而写作
\begin{equation}\label{b2:eq:08-characteristic-differential}
 a(\mathrm dy)^2-2b\,\mathrm dx\,\mathrm dy+c(\mathrm dx)^2=0.
\end{equation}
当 \(\mathrm dx\ne0\) 时，它变为斜率方程 \(am^2-2bm+c=0\)。这解释了为什么这里混合项的符号与主符号 \(a\xi^2+2b\xi\eta+c\eta^2\) 相反：一个写切向，一个写法向。

\ProseParagraphBreak

非特征条件的代数意义是可以解出法向二阶导数。取坐标 \(y_1=\phi(z)\)，补足其余坐标后，\(\partial_{y_1}^2v\) 的系数就是 \(p(z,\nabla\phi)\)。该系数非零时，可以把它移到方程一侧并相除。但这一步只整理了方程，没有证明给定 Cauchy 数据的解存在，也没有得到稳定性。例如对 Laplace 方程，任意非零实法向都是非特征的，而给定函数值和法向导数的 Cauchy 问题仍可能非常不稳定。

\subsection{坐标不变性}
\label{b2:subsec:080104}

\begin{proposition}\label{b2:prop:08-coordinate-change}
设 \(n\geq1\)，\(\Omega,\widetilde\Omega\subset\mathbb R^n\) 为开集，实系数算子 \(L=\sum_{i,j}a^{ij}(z)\partial_{ij}+\sum_i b^i(z)\partial_i+c(z)\) 的系数均连续，且 \(A=(a^{ij})\) 为实对称矩阵场。设 \(F:\Omega\to\widetilde\Omega\) 为 \(C^2\) 微分同胚，\(\mathbb K\in\{\mathbb R,\mathbb C\}\)，\(v\in C^2(\widetilde\Omega;\mathbb K)\)，\(y=F(z)\)，\(u(z)=v(F(z))\)，并令 \(B=DF(z)\)。则新坐标中的主部矩阵为
\begin{equation}\label{b2:eq:08-congruence}
 \widetilde A(y)=B A(z)B^{\mathsf T},\quad z=F^{-1}(y),
\end{equation}
一阶系数为
\begin{equation}\label{b2:eq:08-lower-transform}
 \widetilde b^k(y)=\sum_{i,j}a^{ij}(z)\partial_{ij}F^k(z)
 +\sum_i b^i(z)\partial_iF^k(z).
\end{equation}
零阶系数为 \(\widetilde c(y)=c(F^{-1}(y))\)，主符号满足 \(\widetilde p(y,\eta)=p(z,B^{\mathsf T}\eta)\)。因此，对 \(f\in C(\Omega;\mathbb K)\)，方程 \(Lu=f\) 等价于
\[
 \sum_{k,l}\widetilde a^{kl}\partial_{kl}v
 +\sum_k\widetilde b^k\partial_kv+\widetilde c\,v
 =\widetilde f,
 \qquad \widetilde f=f\circ F^{-1}.
\]
\end{proposition}
\begin{proof}
链式法则给出
\[
 \partial_i u=\sum_k\partial_kv\cdot\partial_iF^k,
\quad
 \partial_{ij}u=\sum_{k,l}\partial_{kl}v\cdot\partial_iF^k\cdot\partial_jF^l
 +\sum_k\partial_kv\cdot\partial_{ij}F^k.
\]
代入算子并按 \(v\) 的导数阶数整理，得到两条系数公式。再把 \(\eta\) 代入二次型，便有
\(\eta^{\mathsf T}BAB^{\mathsf T}\eta=(B^{\mathsf T}\eta)^{\mathsf T}A(B^{\mathsf T}\eta)\)。
\end{proof}

实对称矩阵在可逆\term[hetong-bianhuan]{合同变换}[congruence transformation]下保持正、负及零特征值的个数，这三项称为其\term[guanxing]{惯性}[inertia]。为说明这一事实，先由有限维谱定理把矩阵正交对角化。二次型为正定的子空间最大维数恰等于正特征值个数：正特征子空间达到这个维数；任何更大子空间必与非正特征子空间非平凡相交，不可能正定。可逆线性变换保持子空间维数，因而保持该最大值。对负号二次型重复此论证可得负特征值个数不变，最后由总维数得到零特征值个数不变。

\ProseParagraphBreak

所以主部分类不依赖合法坐标，但矩阵的具体特征值和统一常数可能改变。设 \(A(z)\geq\lambda I\)，其中 \(\lambda>0\) 与 \(z\) 无关，则
\[
 \eta^{\mathsf T}\widetilde A(F(z))\eta
 \geq\lambda|B(z)^{\mathsf T}\eta|^2.
\]
若还知道 \(\sup_{z\in\Omega}\|B(z)^{-1}\|\leq M<\infty\)，由
\( |\eta|\leq\|B(z)^{-1}\|\,|B(z)^{\mathsf T}\eta|\)
得到 \(\widetilde A\geq\lambda M^{-2}I\)。这里均使用欧氏算子范数，且转置不改变该范数。逆导数的一致控制是从 \(A\geq\lambda I\) 这一已知下界推出新下界的充分条件；缺少它时，这一个下界本身不足以保证结论，仍须检查 \(A\) 与 \(B\) 的共同变化。

\ProseParagraphBreak

例如在 \(\mathbb R^2\) 上取
\[
 A(x,y)=\begin{pmatrix}1&0\\0&(1+y^2)^2\end{pmatrix},
 \qquad F(x,y)=(x,\arctan y).
\]
这是到 \(\mathbb R\times(-\pi/2,\pi/2)\) 的光滑微分同胚，且
\[
 B=\begin{pmatrix}1&0\\0&(1+y^2)^{-1}\end{pmatrix},
 \qquad \|B^{-1}\|=1+y^2,
 \qquad BAB^{\mathsf T}=I.
\]
逆导数虽然无界，原矩阵在同一方向上的增长却恰好抵消坐标压缩，新矩阵仍有统一正下界。反过来，若保持同一个 \(F\) 而取 \(A=I\)，新矩阵第二个特征值为 \((1+y^2)^{-2}\)，确实没有统一正下界。这两个计算同时说明充分条件的用途和它的适用边界。

\ProseParagraphBreak

还须区分统一正下界与\secref{b2:sec:0802}采用的双侧一致椭圆条件。前一个例子的原矩阵没有统一上界，不能称为该双侧意义下的一致椭圆矩阵。若已知 \(\lambda I\leq A\leq\Lambda I\)，再有 \(\sup\|B\|\leq N\)、\(\sup\|B^{-1}\|\leq M\)，则
\[
 \frac{\lambda}{M^2}|\eta|^2
 \leq\eta^{\mathsf T}\widetilde A\eta
 \leq\Lambda N^2|\eta|^2,
\]
从而双侧条件也被传递。点态坐标不变性与全域定量估计是两个层次。非线性坐标还会产生\eqref{b2:eq:08-lower-transform}中的一阶项，忽略它们会改变原方程。

% !TeX root = ../../Book2.tex
\section{Gauss 函数与自对偶性}
\label{b2:sec:0602}

Poisson 求和把一个函数的格点和转成其变换的格点和。Gauss 函数的变换仍为 Gauss 函数，因此可以显式比较大尺度和小尺度，并把这种比较放回热传播。

\subsection{Gauss 函数的 Fourier 变换}
\label{b2:subsec:060201}

记 \(g_t(x)=e^{-\pi t x^2}\)，其中 \(t>0\)。由\cref{b2:thm:03-gaussian}及伸缩公式，
\[
 \widehat {g_t}(\xi)=t^{-1/2}e^{-\pi\xi^2/t}.
\]
与保持单位质量的 \(G_\varepsilon\) 相比，\(g_t\) 在原点的高度固定为 \(1\)，总积分为 \(t^{-1/2}\)。将两类参数混用会漏掉平方根因子；这里分别用 \(t\) 和 \(\varepsilon\) 保留区别。

\ProseParagraphBreak

平移和调制给出更一般的波包公式。对实数 \(a,b\)，
\[
 \mathcal F\bigl(e^{-\pi t(x-a)^2}e^{2\pi\mathrm{i}bx}\bigr)(\xi)
 =t^{-1/2}e^{-2\pi\mathrm{i}a(\xi-b)}e^{-\pi(\xi-b)^2/t}.
\]
因此空间中心 \(a\) 产生相位，调制频率 \(b\) 产生频率中心，而宽度参数倒置。每一项都来自\chapref{b2:ch:03}已经证明的运算规则。

\subsection{自对偶结构}
\label{b2:subsec:060202}

对本卷 Fourier 变换已定义的函数或分布 \(f\)，若 \(\widehat f=f\) 在同一空间中成立，称 \(f\) 对当前 Fourier 变换\term[zi-duiou]{自对偶}[self-dual]。Gauss 函数 \(g_1\) 满足这一条件。一般 \(g_t\) 的参数变为 \(1/t\)，并附带 \(t^{-1/2}\)；只有 \(t=1\) 保持同一个归一化函数。

\ProseParagraphBreak

\subsecref{b2:subsec:050204}的 Hermite 函数提供了正 Gauss 函数之外的自对偶例子。\cref{b2:thm:05-hermite}给出 \(\widehat h_n=(-\mathrm i)^nh_n\)，故对每个 \(j\in\mathbb N_0\)，\(h_{4j}\) 都自对偶，有限复线性组合仍然自对偶。相反，\(h_{4j+2}\) 变换后变号，奇数指标对应虚特征值。自对偶是一项线性变换关系，不意味着非负性或概率密度性质。

\subsection{热核解释}
\label{b2:subsec:060203}

\chapref{b2:ch:05}已使用 \(H_t(x)=(4\pi t)^{-1/2}e^{-x^2/(4t)}\)，其变换为 \(e^{-4\pi^2t\xi^2}\)。在实线上周期化，得到
\begin{equation}\label{b2:eq:06-periodic-heat}
 K_t(x)\coloneqq\sum_{n\in\mathbb Z}H_t(x+n)
 =\sum_{k\in\mathbb Z}e^{-4\pi^2tk^2}e^{2\pi\mathrm{i}kx}.
\end{equation}
第二个表达由 Poisson 求和得到；对固定 \(t>0\)，两种表达都绝对收敛。函数 \(K_t\) 称为圆周上的\term[rehe]{热核}[heat kernel]。
\symidx{\(K_t\)：圆周热核}

\begin{proposition}\label{b2:prop:06-heat-kernel}
对 \(t,s>0\)，令 \(K_t(x)=\sum_{n\in\mathbb Z}(4\pi t)^{-1/2}e^{-(x+n)^2/(4t)}\)，在 \(\mathbb T=\mathbb R/\mathbb Z\) 上取总质量为 \(1\) 的 Lebesgue 测度，\(*\) 表示圆周卷积。圆周热核 \(K_t\) 非负、积分为 \(1\)，满足 \(K_t*K_s=K_{t+s}\) 和 \(\partial_tK_t=\partial_x^2K_t\)。对 \(f\in L^p(\mathbb T;\mathbb C)\)、\(1\leq p<\infty\)，当 \(t\downarrow0\) 时有 \(K_t*f\to f\) 于 \(L^p\)；对 \(f\in C(\mathbb T;\mathbb C)\) 则一致收敛。
\end{proposition}
\begin{proof}
非负性由实线表达给出，非负项的 Tonelli 换序将其圆周积分化为实线 Gauss 积分，故积分为 \(1\)。两个核均可积，圆周卷积的 Fourier 系数等于两核系数之积，即 \(e^{-4\pi^2(t+s)k^2}\)；由\cref{b2:cor:02-uniqueness}的周期系数唯一性及两侧连续性得到半群关系。对 \(t\geq\varepsilon>0\)，任意阶 \(x,t\) 导数只在系数上增加一个固定次幂的 \(k\)，仍被可和的指数尾统一控制，故可逐项求导，并得到热方程。

为证初值逼近，对 \(0<\delta<1/2\)，圆周上离原点至少为 \(\delta\) 的区域在周期化前对应到实线中 \(|y|\geq\delta\) 的子集。因此
\[
 \int_{\mathrm{d}_{\mathbb T}(x,0)\geq\delta}K_t(x)\,\mathrm{d}x
 \leq\int_{|y|\geq\delta}H_t(y)\,\mathrm{d}y\longrightarrow0.
\]
所以它满足\cref{b2:prop:02-positive-approximation}中的三项条件：非负、单位质量以及远离原点的质量趋零。该命题对任意趋零的正时间列都适用，因而得到 \(t\downarrow0\) 时的全部收敛结论。
\end{proof}

小时间时，实线表达用少量邻近格点就能描述局部尖峰；大时间时，频率表达只需保留少量低频，就能看出函数趋向常数 \(1\)。两个公式描述同一个核，却在相反的尺度下各自简洁。下一节的 \(\theta\) 变换正是这种尺度对换在 \(x=0\) 处的表达。

\ProseParagraphBreak

这种比较可以给出两端都可核算的误差。对 \(t\geq t_0>0\)，由 \(k^2\geq1+3(k-1)\)、\(k\geq1\)，频率表达给出
\[
 \|K_t-1\|_\infty
 \leq2\sum_{k=1}^\infty e^{-4\pi^2tk^2}
 \leq\frac{2e^{-4\pi^2t}}{1-e^{-12\pi^2t}}.
\]
因此向均匀分布趋近的首个指数由 \(k=\pm1\) 决定。对 \(f\in L^2(\mathbb T;\mathbb C)\)，记平均值 \(\overline f=\int_0^1f\)，这里上横线表示平均而非复共轭。\cref{b2:cor:02-parseval}将 \(L^2\) 范数平方写成全部 Fourier 系数模平方之和，给出精确的能量表达
\[
 \|K_t*f-\overline f\|_2^2
 =\sum_{k\ne0}e^{-8\pi^2tk^2}|\widehat f(k)|^2
 \leq e^{-8\pi^2t}\|f-\overline f\|_2^2.
\]
取 \(f(x)=e^{2\pi\mathrm{i}x}\) 时达到等号，所以能量范数中的衰减率 \(4\pi^2\) 不能统一增大。常数模态保留平均，其余模态耗散。

\ProseParagraphBreak

对于小时间，固定 \(0<a<1/2\) 并取 \(|x|\leq a\)。零格点项为 \(H_t(x)\)，其余项满足 \(|x+n|\geq|n|-a\)。对 \(n\geq1\)，
\[
 (n-a)^2\geq(1-a)^2+(3-2a)(n-1),
\]
故
\[
 0\leq K_t(x)-H_t(x)
 \leq\frac{2(4\pi t)^{-1/2}e^{-(1-a)^2/(4t)}}
              {1-e^{-(3-2a)/(4t)}}.
\]
当 \(0<t\leq1\) 时，分母有依赖 \(a\) 的统一正下界。于是邻近原点的圆周传播，与实线 Gauss 传播之间只有指数小的误差。若接近半周期端点，就可能有两个相邻格点同样重要；单独保留零格点的近似范围必须随之改变。

% !TeX root = ../../Book2.tex
\section{坐标变换}
\label{b2:sec:0803}

主符号的惯性告诉我们可以期待哪一种模型，坐标变换则把这一判断落实为计算。本节区分三种任务：常系数主部在整个坐标域中的标准化，变系数主部在一点的标准化，以及利用特征曲线在一个邻域中构造坐标。三者使用的条件不同，结论也不相同。

\subsection{二维标准型}
\label{b2:subsec:080301}

\begin{theorem}\label{b2:thm:08-constant-canonical}
设 \(a,b,c\) 为实常数且不全为零。存在可逆实线性坐标变换，并允许把整个方程乘以一个非零常数，使主部 \(a\partial_x^2+2b\partial_x\partial_y+c\partial_y^2\) 依照 \(D=b^2-ac\) 的符号分别成为
\[
 \partial_r^2+\partial_s^2\quad(D<0),\quad
 \partial_r^2-\partial_s^2\quad(D>0),\quad
 \partial_r^2\quad(D=0).
\]
低阶项须按同一个变换保留。双曲标准型还可用 \(\xi=r+s\)、\(\eta=r-s\) 化为 \(4\partial_\xi\partial_\eta\)。
\end{theorem}
\begin{proof}
先设 \(a\ne0\)，取 \(r=x\)、\(s=y-(b/a)x\)。此变换行列式为一，且
\[
 \partial_x=\partial_r-\frac ba\partial_s,\quad
 \partial_y=\partial_s.
\]
代入并展开，混合项的系数 \(-2b+2b\) 抵消，得到
\[
 a\partial_r^2+\delta\partial_s^2,\quad
 \delta=c-\frac{b^2}{a}=-\frac Da.
\]
若 \(\delta\ne0\)，再令 \(R=r/\sqrt{|a|}\)、\(S=s/\sqrt{|\delta|}\)，主部变为
\(\operatorname{sgn}(a)\partial_R^2+\operatorname{sgn}(\delta)\partial_S^2\)。乘以 \(\operatorname{sgn}(a)\)，第二项符号等于 \(-\operatorname{sgn}(D)\)，得到前两种形式。若 \(\delta=0\)，只缩放 \(r\) 即得第三种。

若 \(a=0,c\ne0\)，交换 \(x,y\) 后归入上一情形。若 \(a=c=0,b\ne0\)，取 \(r=(x+y)/\sqrt2\)、\(s=(x-y)/\sqrt2\)，则 \(2b\partial_x\partial_y=b(\partial_r^2-\partial_s^2)\)，再缩放并调整符号即可。最后，\(\partial_r=\partial_\xi+\partial_\eta\)、\(\partial_s=\partial_\xi-\partial_\eta\)，相减平方得到 \(4\partial_\xi\partial_\eta\)。
\end{proof}

例如主部 \(4u_{xx}+4u_{xy}+2u_{yy}\) 取 \(r=x/2\)、\(s=y-x/2\) 后成为 \(v_{rr}+v_{ss}\)。同样的变换把
\[
 4u_{xx}+4u_{xy}+u_{yy}+u_x=0
\]
化为 \(v_{rr}+v_r/2-v_s/2=0\)，即 \(v_s=2v_{rr}+v_r\)。二阶部分秩为一，而原来的一阶项恰好提供了新的时间导数。若删掉原式的 \(u_x\)，新方程只是 \(v_{rr}=0\)，便没有这一演化含义。

\ProseParagraphBreak

上述论证适用于常系数。若每一点分别对 \(A(x,y)\) 选一个对角化矩阵，并不能立即把这些矩阵拼成坐标变换的导数：一个候选矩阵场必须满足混合偏导的相容条件。特别是把位置依赖的向量场平方，会额外产生系数导数；不能把变系数算子当成常系数多项式直接配方而忽略这些项。

\subsection{特征坐标}
\label{b2:subsec:080302}

先处理一个双曲邻域。假定 \(a,b,c\in C^2\)、\(a\ne0\)、\(D=b^2-ac>0\)。二维特征斜率为
\[
 m_\pm=\frac{b\pm\sqrt D}{a}.
\]
两族曲线分别由 \(y'=m_\pm(x,y)\) 给出。\chapref{b2:ch:07}的局部流与初值可微依赖保证：以横截直线 \(x=x_0\) 上的交点作为标签，可构造局部函数 \(\xi,\eta\)，使
\[
 \xi_x+m_+\xi_y=0,\quad
 \eta_x+m_-\eta_y=0,\quad \xi_y\ne0,\ \eta_y\ne0.
\]
它们分别在两族特征线上保持常值，称为\term[tezheng-zuobiao]{特征坐标}[characteristic coordinates]。系数为 \(C^2\) 时，流也是 \(C^2\)：对其一阶变分方程再求一次初值导数，所得二阶导数满足线性方程，其非齐次项为速度的二阶导数作用于两个一阶变分；局部一致连续性控制差商余项，再用 Gronwall 不等式得到导数存在且连续。因此这里的坐标变换符合二阶链式法则的要求。

\begin{proposition}\label{b2:prop:08-hyperbolic-coordinates}
上述 \((\xi,\eta)\) 在充分小的邻域内构成 \(C^2\) 微分同胚，并把方程化为
\[
 \gamma(\xi,\eta)v_{\xi\eta}
 +\beta_1(\xi,\eta)v_\xi+\beta_2(\xi,\eta)v_\eta
 +\kappa(\xi,\eta)v=\widetilde f,
 \quad \gamma\ne0.
\]
\end{proposition}
\begin{proof}
坐标的行列式为
\[
 \xi_x\eta_y-\xi_y\eta_x
 =(m_--m_+)\xi_y\eta_y\ne0,
\]
因此可用逆函数定理。由特征方程，\(p(\nabla\xi)=p(\nabla\eta)=0\)，新矩阵的两个对角元为零。其混合系数为
\[
 \gamma=2\bigl(a\xi_x\eta_x+b(\xi_x\eta_y+\xi_y\eta_x)+c\xi_y\eta_y\bigr).
\]
合同变换保持矩阵可逆，而 \(\left(\begin{smallmatrix}0&\gamma/2\\\gamma/2&0\end{smallmatrix}\right)\) 可逆恰好要求 \(\gamma\ne0\)。一阶与零阶系数按\cref{b2:prop:08-coordinate-change}保留，即得所述方程。
\end{proof}

例如在 \(x>0\) 上，方程 \(u_{xx}-x^2u_{yy}=0\) 的两族特征为 \(y\pm x^2/2\) 恒定。取 \(\xi=y+x^2/2\)、\(\eta=y-x^2/2\)，其行列式为 \(2x\ne0\)。两次求导得到
\[
 u_{xx}-x^2u_{yy}=-4x^2v_{\xi\eta}+v_\xi-v_\eta.
\]
新方程不只是 \(v_{\xi\eta}=0\)；两个一阶项正是坐标曲率带来的贡献。在 \(x=0\) 上行列式退化，所以不能用同一坐标穿过那里。

\ProseParagraphBreak

若 \(D=0\) 在整个邻域成立，且 \(a\ne0\)，令 \(\beta=b/a\)、\(V=\partial_x+\beta\partial_y\)，则
\[
 a\partial_x^2+2b\partial_x\partial_y+c\partial_y^2
 =aV^2-a(\partial_x\beta+\beta\partial_y\beta)\partial_y.
\]
取 \(\xi=x\)，并沿 \(V\) 的流构造 \(\eta\) 使 \(V\eta=0\)、\(\eta_y\ne0\)，则 \(V=\partial_\xi\)。方程成为
\[
 a\,v_{\xi\xi}+\beta_1v_\xi+\beta_2v_\eta+\kappa v=\widetilde f.
\]
这是邻域中的秩一标准形式。只有 \(\beta_2\ne0\) 时才能解出 \(v_\eta\)，要按 \(\eta\) 正方向解释为向前热型，还须 \(-a/\beta_2>0\)。秩为一的二阶部分本身没有决定这些条件。

\ProseParagraphBreak

椭圆型没有两族实特征线，因此上述沿实轨道找坐标的方法不适用。一般变系数椭圆主部的邻域标准化需要另解坐标方程，本章不把它当作已证工具；本节已经完整建立的是常系数标准型及一点处冻结系数的标准型。某些特殊系数可直接处理，例如在全平面上取 \(r=e^x\)、\(s=y\)，有
\[
 u_{xx}+e^{2x}u_{yy}=r^2(v_{rr}+v_{ss})+r v_r.
\]
这个映射是全平面到半平面的微分同胚。除以 \(r^2\) 后主部化为 Laplace 算子，但保留了 \(r^{-1}v_r\)；所获表达式的定义域也随坐标改变。

\subsection{高维主符号}
\label{b2:subsec:080303}

在 \(n\) 维，实对称矩阵可经正交变换化为 \(\operatorname{diag}(\lambda_1,\ldots,\lambda_n)\)，再沿非零特征方向缩放，使非零对角元都成为 \(1\) 或 \(-1\)。因此常系数主部完全由惯性三元组 \((n_+,n_-,n_0)\) 决定，允许整体乘以负数时还可交换前两项。这个结论包含二维的三种标准型，但高维中还有更多可能。

\ProseParagraphBreak

例如 \(\partial_1^2+\partial_2^2-\partial_3^2-\partial_4^2\) 的惯性为 \((2,2,0)\)。它既不椭圆，也不属于只有一个异号方向的双曲类型，更不是抛物型。又如三变量算子 \(\partial_1^2\) 的惯性为 \((1,0,2)\)，有两个完全缺失的二阶方向。对这些算子只说“判别式非正”没有明确含义；一个行列式也不能恢复全部惯性信息。

\ProseParagraphBreak

变系数时，可以在任意固定点 \(z_0\) 选一个线性坐标使 \(A(z_0)\) 标准化。邻近点的矩阵通常只是接近标准矩阵。若该点主部非退化，连续性使惯性在足够小邻域内保持不变；理由是单位球上可逆矩阵的最小伸缩量有正下界，小扰动不可能使某个特征值穿过零。秩退化情形没有这样的稳定性，任意小的系数扰动都可能改变类型。

\subsection{几何意义}
\label{b2:subsec:080304}

主符号作用在法向上，特征曲线却画在自变量空间里。两者之间的联系可以用一个直接计算辨明。若 \(\nu\ne0\)、\(p(z,\nu)=0\)，则
\[
 \nu\cdot A(z)\nu=0.
\]
所以当 \(A\nu\ne0\) 时，向量 \(A\nu\) 位于法向为 \(\nu\) 的切平面中。对于非退化双曲型，\(A\nu\) 必非零；对于秩退化类型，它可能恰为零。二阶主符号不仅记录符号，也记录哪些方向上二阶约束消失。

\ProseParagraphBreak

在一维波动方程中，特征曲线为 \(x\pm ct\) 恒定，两族曲线把初值上的区间与后来一点联系起来，\chapref{b2:ch:09}将据此推导表示公式。高维波动方程对应实特征锥，\chapref{b2:ch:11}进一步研究其传播区域。热方程的时间方向由一阶项选定，不能把秩退化的零方向画成波动锥；Laplace 方程则没有非零实特征法向。坐标标准化明确了这些差别，具体解的性质仍由后续各章在相应模型中建立。

% !TeX root = ../../Book2.tex
\OptionalSection{高维 Poisson 求和}{b2:sec:0604}

整数格的直角坐标并非本质要求。任意全秩格都可以由可逆线性映射得到，频率侧相应出现逆转置映射。除了点的位置，还必须追踪基本区域的体积。

\subsection{格与对偶格}
\label{b2:subsec:060401}

\begin{definition}\label{b2:def:06-lattice}
设 \(d\geq1\) 为整数。若 \(A\) 是可逆实 \(d\times d\) 矩阵，则集合 \(\Lambda=A\mathbb Z^d\) 称为 \(\mathbb R^d\) 中的\term[ge]{格}[lattice]，本节均指全秩格。集合 \(A[0,1)^d\) 为一个基本区域，其体积
\[
 v(\Lambda)\coloneqq|\det A|
\]
称为格的\term[jiben-quyu-tiji]{基本区域体积}[covolume]。定义\term[duiou-ge]{对偶格}[dual lattice]为
\[
 \Lambda^*\coloneqq\{\,\xi\in\mathbb R^d:\xi\cdot v\in\mathbb Z\quad(\forall v\in\Lambda)\,\}.
\]
\end{definition}
\symidx{\(\Lambda^*\)：对偶格}
\symidx{\(v(\Lambda)\)：格的基本区域体积}

\begin{proposition}\label{b2:prop:06-dual-lattice}
设 \(d\geq1\) 为整数，\(A\) 是可逆实 \(d\times d\) 矩阵，\(\Lambda=A\mathbb Z^d\)，\(\Lambda^*=\{\,\xi\in\mathbb R^d:\xi\cdot v\in\mathbb Z\ (v\in\Lambda)\,\}\)，\(v(\Lambda)=|\det A|\)。则 \(\Lambda^*=A^{-\mathsf T}\mathbb Z^d\)、\((\Lambda^*)^*=\Lambda\) 和 \(v(\Lambda^*)=v(\Lambda)^{-1}\)。基本区域体积不依赖格基的选择。
\end{proposition}
\begin{proof}
条件 \(\xi\cdot An\in\mathbb Z\) 对全部整数向量 \(n\) 成立，等价于 \(A^{\mathsf T}\xi\) 的每个坐标都是整数，所以得到对偶格公式。再次取逆转置和行列式便得其余两个等式。若另一基矩阵为 \(B\)，则 \(U=A^{-1}B\) 及 \(U^{-1}\) 都是整数矩阵，因此其行列式为互为倒数的整数，只能为 \(\pm1\)。故 \(|\det B|=|\det A|\)。
\end{proof}

对角矩阵给出矩形格，边长分别为正数 \(a_j\) 时，对偶边长为 \(a_j^{-1}\)。非正交基也遵循逆转置，而不是把每个基向量分别除以其长度；内积的整数条件才是定义。

\Cref{b2:fig:06-dual-lattice}用一个非正交格显示逆转置的作用。箭头端点的内积满足 \(b_i\cdot b_j^*=\delta_{ij}\)，不能用逐个基向量取倒长度代替这一条件。
% 本书原创：同一欧氏比例绘出非正交格与逆转置对偶格。
\begin{figure}[htbp]
\centering
\begin{tikzpicture}[x=1cm,y=1cm,>=stealth,every node/.style={font=\small},line width=.65pt]
\begin{scope}[xshift=0.0cm]
\fill[AnalysisBlue!10] (0.00000,0.00000)--(1.50000,0.00000)--(2.00000,1.00000)--(0.50000,1.00000)--cycle;
\draw[->,gray] (-1.2,0)--(3.4,0);\draw[->,gray] (0,-1.2)--(0,2.6);
\fill[gray] (-1.00000,1.00000) circle(1.15pt);
\fill[gray] (-0.50000,2.00000) circle(1.15pt);
\fill[gray] (-0.50000,-1.00000) circle(1.15pt);
\fill[gray] (0.00000,0.00000) circle(1.15pt);
\fill[gray] (0.50000,1.00000) circle(1.15pt);
\fill[gray] (1.00000,2.00000) circle(1.15pt);
\fill[gray] (1.00000,-1.00000) circle(1.15pt);
\fill[gray] (1.50000,0.00000) circle(1.15pt);
\fill[gray] (2.00000,1.00000) circle(1.15pt);
\fill[gray] (2.50000,2.00000) circle(1.15pt);
\fill[gray] (2.50000,-1.00000) circle(1.15pt);
\fill[gray] (3.00000,0.00000) circle(1.15pt);
\draw[->,AnalysisBlue,thick] (0,0)--(1.50000,0.00000) node[above right] {\(b_1\)};
\draw[->,AnalysisBlue,thick] (0,0)--(0.50000,1.00000) node[above right] {\(b_2\)};
\node at (1,3.0) {\(\Lambda,\quad v(\Lambda)=3/2\)};\end{scope}
\begin{scope}[xshift=6.6cm]
\fill[AnalysisBlue!10] (0.00000,0.00000)--(0.66667,-0.33333)--(0.66667,0.66667)--(0.00000,1.00000)--cycle;
\draw[->,gray] (-1.2,0)--(3.4,0);\draw[->,gray] (0,-1.2)--(0,2.6);
\fill[gray] (-0.66667,-0.66667) circle(1.15pt);
\fill[gray] (-0.66667,0.33333) circle(1.15pt);
\fill[gray] (-0.66667,1.33333) circle(1.15pt);
\fill[gray] (0.00000,-1.00000) circle(1.15pt);
\fill[gray] (0.00000,0.00000) circle(1.15pt);
\fill[gray] (0.00000,1.00000) circle(1.15pt);
\fill[gray] (0.00000,2.00000) circle(1.15pt);
\fill[gray] (0.66667,-0.33333) circle(1.15pt);
\fill[gray] (0.66667,0.66667) circle(1.15pt);
\fill[gray] (0.66667,1.66667) circle(1.15pt);
\fill[gray] (1.33333,-0.66667) circle(1.15pt);
\fill[gray] (1.33333,0.33333) circle(1.15pt);
\fill[gray] (1.33333,1.33333) circle(1.15pt);
\fill[gray] (2.00000,-1.00000) circle(1.15pt);
\fill[gray] (2.00000,0.00000) circle(1.15pt);
\fill[gray] (2.00000,1.00000) circle(1.15pt);
\fill[gray] (2.00000,2.00000) circle(1.15pt);
\fill[gray] (2.66667,-0.33333) circle(1.15pt);
\fill[gray] (2.66667,0.66667) circle(1.15pt);
\fill[gray] (2.66667,1.66667) circle(1.15pt);
\draw[->,AnalysisBlue,thick] (0,0)--(0.66667,-0.33333) node[above right] {\(b_1^*\)};
\draw[->,AnalysisBlue,thick] (0,0)--(0.00000,1.00000) node[above right] {\(b_2^*\)};
\node at (1,3.0) {\(\Lambda^*,\quad v(\Lambda^*)=2/3\)};\end{scope}
\end{tikzpicture}
\caption[非正交格与对偶格]{非正交格与对偶格。左侧基矩阵为 \(\left(\begin{smallmatrix}3/2&1/2\\0&1\end{smallmatrix}\right)\)，右侧取逆转置；淡蓝色区域为相应基本平行四边形。两图采用同一长度单位，区域面积互为倒数。}
\label{b2:fig:06-dual-lattice}
\end{figure}


\subsection{欧氏空间中的求和公式}
\label{b2:subsec:060402}

\begin{theorem}[全秩格上的 Poisson 求和]\label{b2:thm:06-lattice-poisson}
设 \(d\geq1\) 为整数，\(f\in\mathcal S(\mathbb R^d;\mathbb C)\)，\(\Lambda=A\mathbb Z^d\)，其中 \(A\) 是可逆实 \(d\times d\) 矩阵。令 \(v(\Lambda)=|\det A|\)，\(\Lambda^*=A^{-\mathsf T}\mathbb Z^d\)。则
\[
 \sum_{v\in\Lambda}f(v)
 =\dfrac1{v(\Lambda)}\sum_{\xi\in\Lambda^*}\widehat f(\xi).
\]
两侧都绝对收敛。
\end{theorem}
\begin{proof}
先取 \(\Lambda=\mathbb Z^d\)。逐坐标应用一维 Poisson 公式；固定其他坐标时，一维切片为 Schwartz 函数。为核对每步的共同控制，先看第一个坐标的部分变换
\[
 g(\eta,x')=\int_{\mathbb R}f(x_1,x')e^{-2\pi\mathrm{i}x_1\eta}\,\mathrm{d}x_1.
\]
对 \(\eta\) 求导在积分内产生 \(x_1\) 的幂；乘以 \(\eta\) 的幂可改成 \(x_1\) 方向的导数，再分部积分；对 \(x'\) 的导数与权重则直接保留。这样每个目标加权导数都是某个 \(f\) 的加权导数在 \(x_1\) 上的积分。用足够高的原半范数，可使被积量由 \(C(1+|x_1|)^{-2}\) 控制，同时保留任意指定阶的 \(\eta,x'\) 权重。因此 \(g\) 仍为全部 \(d\) 个变量上的 Schwartz 函数。后续坐标变换重复同一论证。

在每一个中间阶段，进一步取足够高的权重，便可用 \(C\prod_j(1+|y_j|)^{-2}\) 控制混合变量函数；\(y_j\) 可以是待积分的实变量，也可以是已取样的整数指标。连续积分和离散级数均可和，故 Fubini 与绝对收敛重排适用。每次一维求和只把一个坐标的整数采样移到该坐标的频率侧，依次得到 \(\sum_{n\in\mathbb Z^d}f(n)=\sum_{k\in\mathbb Z^d}\widehat f(k)\)。

一般情形令 \(g(x)=f(Ax)\)，则 \(g\in\mathcal S\)，且 \(\widehat g(k)=|\det A|^{-1}\cdot\widehat f(A^{-\mathsf T}k)\)。对 \(g\) 使用整数格公式，即得结论。绝对收敛由可逆映射的长度比较及任意高阶的快速下降保证。
\end{proof}

基本区域体积因子不能省略。例如在一维取 \(h>0\)、\(\Lambda=h\mathbb Z\)，对偶格为 \(h^{-1}\mathbb Z\)，公式恰好恢复 \(\sum_nf(hn)=h^{-1}\sum_k\widehat f(k/h)\)。高维中的行列式就是这个因子的体积版本。

\subsection{格的 \texorpdfstring{\(\theta\)}{\AnalysisPdfTheta} 级数}
\label{b2:subsec:060403}

\begin{definition}\label{b2:def:06-lattice-theta}
设 \(d\geq1\) 为整数，\(\Lambda\subset\mathbb R^d\) 为全秩格。对 \(t>0\)，定义它的\term[ge-theta-jishu]{格 \(\theta\) 级数}[lattice theta series]为
\[
 \Theta_\Lambda(t)\coloneqq\sum_{v\in\Lambda}e^{-\pi t|v|^2}.
\]
\end{definition}
\symidx{\(\Theta_\Lambda\)：全秩格的 \(\theta\) 级数}

\begin{corollary}\label{b2:cor:06-lattice-theta}
设 \(d\geq1\) 为整数，\(\Lambda=A\mathbb Z^d\)，其中 \(A\) 是可逆实 \(d\times d\) 矩阵，\(\Lambda^*=A^{-\mathsf T}\mathbb Z^d\) 为其对偶格，\(v(\Lambda)=|\det A|\)。对全秩格 \(L\subset\mathbb R^d\)，记 \(\Theta_L(u)=\sum_{v\in L}e^{-\pi u|v|^2}\)（\(u>0\)）。这些级数收敛，且对每个 \(t>0\)，
\[
 \Theta_\Lambda(t)=v(\Lambda)^{-1}t^{-d/2}\Theta_{\Lambda^*}(1/t).
\]
\end{corollary}
\begin{proof}
在\cref{b2:thm:06-lattice-poisson}中取 \(f(x)=e^{-\pi t|x|^2}\)。对 \(t>0\)，这个函数属于 \(\mathcal S(\mathbb R^d)\)，其变换为 \(t^{-d/2}e^{-\pi|\xi|^2/t}\)。该定理保证两侧级数绝对收敛，代入即得。
\end{proof}

若 \(\Lambda=\mathbb Z^d\)，则对 \(t>0\) 有 \(\Theta_\Lambda(t)=\vartheta(t)^d\)。一般的变换公式同时涉及倒数参数、对偶格及基本区域体积。满足 \(\Lambda^*=\Lambda\) 的格称为\term[ziduiou-ge]{自对偶格}[self-dual lattice]；此时由体积互反可知 \(v(\Lambda)=1\)，两侧直接成为同一个格的级数。自对偶是做到这一点的充分条件。

\ProseParagraphBreak

更一般地，设 \(d\geq1\) 为整数，\(\Lambda\subset\mathbb R^d\) 为全秩格，\(c>0\)，\(U\) 是实 \(d\times d\) 正交矩阵，即 \(U^{\mathsf T}U=I_d\)，并假设 \(\Lambda^*=cU\Lambda\)。正交变换保持长度，伸缩将平方长度乘以 \(c^2\)，所以对 \(u>0\)，由格点的一一对应得到
\[
 \Theta_{\Lambda^*}(u)
 =\sum_{v\in\Lambda}e^{-\pi u c^2|Uv|^2}
 =\Theta_\Lambda(c^2u).
\]
体积也受这一假设约束：由 \(|\det U|=1\)，
\[
 v(\Lambda)^{-1}=v(\Lambda^*)=c^dv(\Lambda),
 \quad v(\Lambda)=c^{-d/2}.
\]
代入\cref{b2:cor:06-lattice-theta}，对每个 \(t>0\) 得
\[
 \Theta_\Lambda(t)
 =v(\Lambda)^{-1}t^{-d/2}\Theta_\Lambda(c^2/t)
 =c^{d/2}t^{-d/2}\Theta_\Lambda(c^2/t).
\]
当 \(c=1\) 时，只要对偶格与原格正交等距，就有 \(\Theta_{\Lambda^*}(t)=\Theta_\Lambda(t)\) 和 \(v(\Lambda)=1\)，从而得到同一格的倒数参数公式；不要求两个格作为点集相同。

\ProseParagraphBreak

例如二维矩形格
\[
 \Lambda=2\mathbb Z\times\tfrac12\mathbb Z,
 \quad \Lambda^*=\tfrac12\mathbb Z\times2\mathbb Z
\]
并不相同，因为 \((0,1/2)\in\Lambda\) 而 \((0,1/2)\notin\Lambda^*\)。交换坐标的正交矩阵 \(U=\left(\begin{smallmatrix}0&1\\1&0\end{smallmatrix}\right)\) 将 \(\Lambda\) 映到 \(\Lambda^*\)，且 \(v(\Lambda)=2\cdot(1/2)=1\)。因此虽然这个格不自对偶，仍有
\[
 \Theta_\Lambda(t)=t^{-1}\Theta_\Lambda(1/t)\quad(t>0).
\]

对可逆实 \(d\times d\) 基矩阵 \(A\)，若用 Gram 矩阵 \(Q=A^{\mathsf T}A\) 表示长度，则对每个 \(n\in\mathbb Z^d\)，
\[
 |An|^2=n^{\mathsf T}Qn,
 \quad |A^{-\mathsf T}n|^2=n^{\mathsf T}Q^{-1}n.
\]
整数换基作用在系数 \(n\in\mathbb Z^d\) 上，欧式正交变换作用在格向量 \(An\in\mathbb R^d\) 上。两者作用的对象不同；即使矩阵写法相近，也须区分。这为第三卷的二次型 \(\theta\) 级数提供几何来源，但不要求在这里提前引入整数性或模形式的附加结构。

\ProseParagraphBreak

作为非直角的可算例子，取
\[
 A=\begin{pmatrix}\sqrt2&1/\sqrt2\\0&\sqrt{3/2}\end{pmatrix},
 \quad Q=A^{\mathsf T}A=\begin{pmatrix}2&1\\1&2\end{pmatrix}.
\]
格基向量的夹角为 \(\pi/3\)，基本区域体积为 \(\sqrt3\)。对 \(t>0\)，记
\[
 B(t)=\sum_{m,n\in\mathbb Z}e^{-2\pi t(m^2+mn+n^2)}.
\]
这是该格的 \(\theta\) 级数。取整数基变换矩阵
\[
 P=\begin{pmatrix}1&0\\0&-1\end{pmatrix}.
\]
它满足 \(P\mathbb Z^2=\mathbb Z^2\)，并且
\[
 Q^{-1}=\frac13\begin{pmatrix}2&-1\\-1&2\end{pmatrix}
 =\frac13P^{\mathsf T}QP.
\]
因此对 \(u>0\)，把整数系数 \(k=(m,n)^{\mathsf T}\) 换成 \(Pk=(m,-n)^{\mathsf T}\)，得到
\[
 \sum_{k\in\mathbb Z^2}e^{-\pi u k^{\mathsf T}Q^{-1}k}
 =\sum_{k\in\mathbb Z^2}e^{-\pi(u/3)(Pk)^{\mathsf T}Q(Pk)}
 =B(u/3).
\]
Poisson 公式于是给出
\[
 B(t)=\frac1{\sqrt3\,t}
 \sum_{m,n\in\mathbb Z}e^{-2\pi(m^2-mn+n^2)/(3t)}
 =\frac1{\sqrt3\,t}B(1/(3t)).
\]
这里参数中的 \(1/3\) 来自二次型的尺度，而 \(P\) 只负责将整数系数重新编号。要在实际欧式空间中写出对应的正交变换，可取
\[
 U_0=\begin{pmatrix}\sqrt3/2&-1/2\\-1/2&-\sqrt3/2\end{pmatrix}.
\]
直接相乘可核对 \(U_0^{\mathsf T}U_0=I_2\) 及 \(A^{-\mathsf T}=3^{-1/2}U_0AP\)。因 \(P\mathbb Z^2=\mathbb Z^2\)，这就说明 \(\Lambda^*=3^{-1/2}U_0\Lambda\)。相似尺度 \(c=1/\sqrt3\) 给出 \(c^2/t=1/(3t)\)，体积关系给出 \(v(\Lambda)=c^{-1}=\sqrt3\)，与上式的两个因子分别相符。

\ProseParagraphBreak

最短非零向量由 \(m^2+mn+n^2=1\) 给出。写成 \((m+n/2)^2+3n^2/4\)，可先知 \(|n|\leq1\)，再逐个求出六组解 \((\pm1,0),(0,\pm1),(1,-1),(-1,1)\)。为明确其余项的统一误差，记 \(q(m,n)=m^2+mn+n^2\)。其余非零向量满足整数条件 \(q(m,n)\geq2\)，并且 \(q(m,n)\geq(m^2+n^2)/2\)。故对 \(t\geq1\)，
\[
 \sum_{q(m,n)\geq2}e^{-2\pi t q(m,n)}
 \leq e^{-4\pi t}\sum_{q(m,n)\geq2}e^{-2\pi(q(m,n)-2)}.
\]
右端的固定级数由二维 Gauss 级数控制而收敛。因此
\[
 B(t)=1+6e^{-2\pi t}+O(e^{-4\pi t})\quad(t\to\infty).
\]
把这一式用于倒数参数，得到
\[
 B(t)=\frac1{\sqrt3\,t}
 \left(1+6e^{-2\pi/(3t)}+O(e^{-4\pi/(3t)})\right)
 \quad(t\downarrow0).
\]
大参数修正读取最短格向量，小参数修正读取最短对偶格向量；主项读取基本区域体积。这三种几何数据在公式中各有确定位置。


\begin{chaptersummary}
周期化把实线上的函数放到圆周上，其 Fourier 系数就是连续变换在整数频率处的值。Poisson 求和从这个事实与系数唯一性得出；尺度变化产生倒数格与体积因子，不能省去对采样点和代表的正则性要求。

\ProseParagraphBreak

\subsecref{b2:subsec:060104}的数值求积把无穷采样误差与有限窗口的尾部误差分开。\subsecref{b2:subsec:060105}另外利用有限频带上的 \(L^2\) 展开研究精确恢复与混叠；这条选读路线可以略过，不影响周期热核及下一编 PDE 内容的学习。

\ProseParagraphBreak

Gauss 变换把参数取倒数，使 \(\theta\) 级数在大小尺度间交换。相同的对称既解释圆周热核的局部尖峰和低频平衡，也把整数平方和与 Mellin 积分联系起来。第三编将在有限区间上使用特征函数展开，在全空间上使用欧氏变换，并用能量范数和分布方程解释低正则解。
\end{chaptersummary}

\BookExercises

题组围绕一维求和、数值求积、周期化的频率采样与平方和计数展开，选做题使用一般格或参数分析。其中的频率采样通过周期化的 Fourier 系数计算；最后的跨章题把周期化与实线、圆周上的 Gauss 平滑联系起来。所有交换极限与求和的步骤均应说明所用的统一控制。

% 自拟：调制Gauss后得到交错theta，避免从复参数模公式直接摘用。
\Needspace{6\baselineskip}
\begin{exercise}\label{b2:ex:06-alternating-theta}
对 \(t>0\)，证明
\[
 \sum_{n\in\mathbb Z}(-1)^ne^{-\pi tn^2}
 =t^{-1/2}\sum_{k\in\mathbb Z}e^{-\pi(k-1/2)^2/t}>0.
\]
求 \(t\downarrow0\) 时的首项，并给出相对误差的指数阶。
\end{exercise}
\begin{solution}
在带调制的 Poisson 公式\cref{b2:cor:06-scaled-poisson}中取 \(f(x)=e^{-\pi tx^2}\)、\(h=1,a=0,b=1/2\)，得到等式；右侧每项为正，故严格为正。距离 \(1/2\) 最近的整数为 \(0,1\)，贡献 \(2t^{-1/2}e^{-\pi/(4t)}\)。将全部半整数距离成对合并，原和式可写成
\[
 2t^{-1/2}e^{-\pi/(4t)}
 \left(1+\sum_{j=1}^\infty e^{-\pi j(j+1)/t}\right).
\]
因为 \(j(j+1)\geq2+4(j-1)\)、\(j\geq1\)，括号中除 \(1\) 外的尾部介于 \(0\) 与 \(e^{-2\pi/t}/(1-e^{-4\pi/t})\) 之间。因此结果为 \(2t^{-1/2}e^{-\pi/(4t)}(1+O(e^{-2\pi/t}))\)。
\end{solution}

% 自拟：theta的自反尺度给出一个精确的带权求和。
\Needspace{6\baselineskip}
\begin{exercise}\label{b2:ex:06-theta-derivative}
证明
\[
 \sum_{n\in\mathbb Z}n^2e^{-\pi n^2}
 =\dfrac1{4\pi}\sum_{n\in\mathbb Z}e^{-\pi n^2}.
\]
\end{exercise}
\begin{solution}
令 \(\vartheta(t)=\sum_{n\in\mathbb Z}e^{-\pi n^2t}\)、\(t>0\)。在 \(t\) 的紧正区间上，导数级数由可和的 Gauss 尾一致控制，故可逐项求导。对\cref{b2:thm:06-theta-transformation}给出的 \(\vartheta(t)=t^{-1/2}\vartheta(1/t)\) 求导并令 \(t=1\)，得 \(\vartheta'(1)=-\vartheta(1)/2-\vartheta'(1)\)，所以 \(\vartheta'(1)=-\vartheta(1)/4\)。再用 \(\vartheta'(1)=-\pi\sum_nn^2e^{-\pi n^2}\) 即得。
\end{solution}

% 自拟：空间截断与频率误差分开控制，可作为实际求积依据。
\Needspace{6\baselineskip}
\begin{exercise}\label{b2:ex:06-quadrature-error}
设 \(h\in(0,1]\)、\(M\) 为非负整数，令 \(Q_{h,M}=h\sum_{|n|\leq M}e^{-\pi h^2n^2}\)。证明
\[
 |Q_{h,M}-1|
 \leq\dfrac{2e^{-\pi/h^2}}{1-e^{-3\pi/h^2}}
 +\dfrac{2h e^{-\pi h^2(M+1)^2}}{1-e^{-\pi h^2(2M+3)}}.
\]
\end{exercise}
\begin{solution}
无穷求和由 Poisson 公式等于 \(\sum_ke^{-\pi k^2/h^2}\)，减去零频项 \(1\) 后，用 \(k^2\geq1+3(k-1)\)、\(k\geq1\) 得到第一个几何级数界。空间截断误差为 \(2h\sum_{n=M+1}^\infty e^{-\pi h^2n^2}\)。写 \(n=M+1+j\)，有 \(n^2\geq(M+1)^2+(2M+3)j\)，得到第二项。两类误差分别来自离散采样和有限截断，最后用三角不等式合并。
\end{solution}

% 改编：Hunter2001Applied, Theorem11.45, p.319 的C1衰减范围；
% 在单位周期下用于非Schwartz有理核，并推导显式双边求和。
\Needspace{6\baselineskip}
\begin{exercise}\label{b2:ex:06-rational-sum}
对 \(a>0\)，以 \(\coth x=(e^x+e^{-x})/(e^x-e^{-x})\)、\(x>0\) 表示本题所需的双曲余切。证明
\[
 \sum_{n\in\mathbb Z}\dfrac1{n^2+a^2}
 =\dfrac\pi a\coth(\pi a).
\]
所用函数不属于 Schwartz 空间，须核对正文的有限阶衰减判据再计算。
\end{exercise}
\begin{solution}
取 \(f(x)=1/(x^2+a^2)\)，其函数和一阶导数分别为 \(O((1+|x|)^{-2})\)、\(O((1+|x|)^{-3})\)。因此满足\cref{b2:prop:06-finite-decay-poisson}的条件，可以使用求和公式。

\subsecref{b2:subsec:030503}的一维 Poisson 核变换说明，\(P_a(x)=a/(\pi(a^2+x^2))\) 满足 \(\widehat{P_a}(\xi)=e^{-2\pi a|\xi|}\)。由 \(f=(\pi/a)P_a\)，得到 \(\widehat f(k)=\pi a^{-1}e^{-2\pi a|k|}\)。在 \(x=0\) 处使用周期化求和公式，得到
\[
 \sum_nf(n)=\dfrac\pi a\left(1+2\sum_{k=1}^\infty e^{-2\pi ak}\right)
 =\dfrac\pi a\dfrac{1+e^{-2\pi a}}{1-e^{-2\pi a}}
 =\dfrac\pi a\coth(\pi a).
\]
\end{solution}

% 自拟：整数采样变换不恢复未周期化的函数。
\Needspace{6\baselineskip}
\begin{exercise}\label{b2:ex:06-periodization-kernel}
构造非零 \(f\in C_c^\infty(\mathbb R;\mathbb C)\)，使周期化 \((\mathcal Pf)(x)=\sum_{n\in\mathbb Z}f(x+n)\) 恒为零。证明其 Fourier 变换在每个整数处为零，但并非处处为零。这是否与变换唯一性矛盾？
\end{exercise}
\begin{solution}
取非零 \(g\in C_c^\infty((-1/4,1/4);\mathbb C)\)，令 \(f(x)=g(x)-g(x-1)\)。两个支集不交，故 \(f\ne0\)，而周期化时两个整数平移和逐项重排后相同，所以 \(\mathcal Pf=0\)。变换为 \((1-e^{-2\pi\mathrm{i}\xi})\widehat g(\xi)\)，在整数处全为零。若它处处为零，\cref{b2:cor:03-uniqueness}会给出 \(f=0\) 几乎处处，连续性再给出 \(f=0\) 处处，矛盾。因此整数频率的采样只确定周期化函数；唯一性定理要求全部连续频率的数据，二者条件不同。
\end{solution}

% 自拟：从平方和生成函数读取有序计数，辨认符号和顺序。
\Needspace{6\baselineskip}
\begin{exercise}\label{b2:ex:06-squares-count}
设 \(d,e\) 为正整数，\(m\) 为非负整数；对每个正整数 \(q\) 和非负整数 \(\ell\)，记 \(r_q(\ell)=\#\{n\in\mathbb Z^q:|n|^2=\ell\}\)，其中符号和坐标顺序均计数。证明 \(r_{d+e}(m)=\sum_{j=0}^mr_d(j)r_e(m-j)\)，并直接计算 \(r_2(5)\)、\(r_4(1)\)、\(r_4(2)\)。解释为何不应再除以排列数。
\end{exercise}
\begin{solution}
把一个 \((d+e)\) 元整数向量分成前 \(d\) 个和后 \(e\) 个坐标，前一组平方和为某个 \(j\)，后一组为 \(m-j\)。对各组独立计数再求和即得公式。平方和为 \(5\) 的二维向量是 \((\pm1,\pm2)\) 及其坐标交换，共 \(8\) 个。四维平方和为 \(1\) 时选一个坐标放 \(\pm1\)，共 \(8\) 个；平方和为 \(2\) 时选两个坐标各放 \(\pm1\)，共 \(\binom42 2^2=24\) 个。theta 乘积的每个因子对应一个有序坐标，因此这些排列已代表不同向量，不应再除掉。
\end{solution}

% 自拟：非正交格的逆转置与最短向量，调用选读正文。
\Needspace{6\baselineskip}
\begin{optionalexercise}\label{b2:ex:06-parity-lattice}
令 \(\Lambda=\{m(2,0)+n(1,1):m,n\in\mathbb Z\}\subset\mathbb R^2\)，并对 \(t>0\) 记 \(\Theta_\Lambda(t)=\sum_{v\in\Lambda}e^{-\pi t|v|^2}\)。求 \(\Lambda^*\)、两格的基本区域体积和最短非零向量长度，写出 \(\theta\) 变换，并求 \(\Theta_\Lambda(t)\) 在 \(t\to\infty\) 时的首个指数修正。
\end{optionalexercise}
\begin{solution}
基矩阵为 \(A=\left(\begin{smallmatrix}2&1\\0&1\end{smallmatrix}\right)\)，故对偶格由 \((1/2,-1/2)\)、\((0,1)\) 生成；基本区域体积分别为 \(2\)、\(1/2\)。原格是两个整数坐标同奇偶的向量：若坐标都为偶数，非零向量长度至少为 \(2\)；若都为奇数，长度至少为 \(\sqrt2\)。因此最短非零向量恰为 \((\pm1,\pm1)\)。对偶格向量可以写成 \((p/2,q-p/2)\)、\(p,q\in\mathbb Z\)。当 \(p\) 为偶数时两坐标都是整数，非零向量长度至少为 \(1\)；当 \(p\) 为奇数时两坐标都是半整数，长度至少为 \(1/\sqrt2\)，取等向量恰为 \((\pm1/2,\pm1/2)\)。因此
\[
 \Theta_\Lambda(t)=\dfrac1{2t}\Theta_{\Lambda^*}(1/t),
 \quad\Theta_\Lambda(t)=1+4e^{-2\pi t}+O(e^{-4\pi t}).
\]
最后的尾部估计可在 \(t\geq1\) 时提出 \(e^{-4\pi t}\)，其余由格上可和的固定 Gauss 级数控制；最短向量后的平方长度至少为 \(4\)。
\end{solution}

% 自拟：正权平方矩的方差解释theta的对数凸性。
\Needspace{6\baselineskip}
\begin{optionalexercise}\label{b2:ex:06-log-convexity}
对 \(t>0\)，令 \(\vartheta(t)=\sum_{n\in\mathbb Z}e^{-\pi n^2t}\)。证明 \(\vartheta\) 在 \((0,\infty)\) 上严格递减，而 \(\log\vartheta\) 在该区间上严格凸。将后一结论解释成以权重 \(w_n(t)=e^{-\pi n^2t}/\vartheta(t)\) 计算的方差为正。
\end{optionalexercise}
\begin{solution}
逐项求导得到 \(\vartheta'(t)=-\pi\sum_nn^2e^{-\pi n^2t}<0\)。权重正且总和为 \(1\)，于是
\[
 (\log\vartheta)''(t)
 =\pi^2\left(\sum_n n^4w_n(t)-\left(\sum_n n^2w_n(t)\right)^2\right).
\]
括号中为随机取值 \(n^2\) 的方差。值 \(0\) 与 \(1\) 都有正权重，所以它不恒定，方差严格为正。所有一、二阶参数求导都由紧正参数区间上的 Gauss 尾控制。
\end{solution}

% 自拟：综合周期系数、实线卷积、Schwartz 衰减和 Poisson 周期化；不以热方程理论为前置。
\Needspace{6\baselineskip}
\begin{exercise}[周期化与 Gauss 平滑的交换]\label{b2:ex:06-periodization-gaussian-commute}
设 \(f\in\mathcal S(\mathbb R;\mathbb C)\)，\(t>0\)，并令
\[
 H_t(x)=(4\pi t)^{-1/2}e^{-x^2/(4t)},\qquad
 K_t=\mathcal PH_t,\qquad
 Q_tf=\mathcal P(H_t*f).
\]
圆周取 \(\mathbb T=\mathbb R/\mathbb Z\) 及总质量为 \(1\) 的 Lebesgue 测度，周期卷积定义为 \((u*_{\mathbb T}v)(x)=\int_0^1u(x-y)v(y)\,\mathrm dy\)。
\begin{enumerate}
\item\label{b2:item:06-gaussian-commute}证明 \(Q_tf=K_t*_{\mathbb T}(\mathcal Pf)\)，并对展开周期化级数后使用的求和、积分交换给出绝对可积控制。
\item\label{b2:item:06-gaussian-coefficients}求 \(Q_tf\) 的全部 Fourier 系数，并证明 \(\int_0^1Q_tf=\int_{\mathbb R}f\)。
\item\label{b2:item:06-gaussian-approximation}对每个非负整数 \(m\)，证明存在只依赖 \(f,m\) 的常数 \(C_{f,m}\)，使 \(0<t\leq1\) 时
\[
 \|Q_tf-\mathcal Pf\|_{C^m(\mathbb T)}\leq C_{f,m}t,
 \qquad
 \|Q_tf-\mathcal Pf-t(\mathcal Pf)''\|_{C^m(\mathbb T)}
 \leq C_{f,m}t^2.
\]
这里 \(\|v\|_{C^m}=\max_{0\leq j\leq m}\|v^{(j)}\|_\infty\)。
\end{enumerate}
\end{exercise}
\begin{solution}
固定 \(t>0\)。由\cref{b2:prop:05-schwartz-convolution}，\(H_t*f\in\mathcal S\)，故其周期化是光滑函数。又 \(K_t\) 是有界的正周期函数。对每个 \(x\)，展开周期卷积时的绝对值总积分满足
\[
 \begin{aligned}
 &\int_0^1\sum_{m,n\in\mathbb Z}|f(y+m)|H_t(x-y+n)\,\mathrm dy\\
 &\quad=\int_{\mathbb R}|f(z)|\sum_{r\in\mathbb Z}H_t(x-z+r)\,\mathrm dz
 \leq\|f\|_1\|K_t\|_\infty<\infty.
 \end{aligned}
\]
其中令 \(z=y+m\)、\(r=n+m\)，各个区间 \([m,m+1)\) 恰好分割实线；对非负项先用 Tonelli 即可得到这个等式。于是对原来的复值积分也可用 Fubini，得到
\[
 \begin{aligned}
 (K_t*_{\mathbb T}\mathcal Pf)(x)
 &=\sum_{r\in\mathbb Z}\int_{\mathbb R}H_t(x-z+r)f(z)\,\mathrm dz\\
 &=\sum_{r\in\mathbb Z}(H_t*f)(x+r)=Q_tf(x).
 \end{aligned}
\]
这说明先把实线输入平滑再周期化，与先周期化再用周期核平滑，得到同一个函数。

\cref{b2:prop:06-periodic-coefficients}和实线卷积定理给出第 \(k\) 个系数
\[
 \widehat{Q_tf}(k)=\widehat{H_t*f}(k)
 =e^{-4\pi^2tk^2}\widehat f(k),\qquad k\in\mathbb Z.
\]
取 \(k=0\) 即得积分恒等式。由于 \(\widehat f\) 仍是 Schwartz 函数，对任意非负整数 \(N\)，级数 \(\sum_k(1+|k|)^N|\widehat f(k)|\) 收敛。因此 Poisson 公式和逐项求导可写为
\[
 Q_tf(x)=\sum_{k\in\mathbb Z}e^{-4\pi^2tk^2}\widehat f(k)e^{2\pi\mathrm ikx},
 \qquad
 \mathcal Pf(x)=\sum_{k\in\mathbb Z}\widehat f(k)e^{2\pi\mathrm ikx}.
\]
对 \(x\geq0\)，有 \(0\leq1-e^{-x}\leq x\) 及 \(|e^{-x}-1+x|\leq x^2/2\)。将这两个不等式分别用于 \(x=4\pi^2tk^2\)，再对 \(j\leq m\) 阶导数逐项估计，所需上界分别为
\[
 4\pi^2t\sum_k(2\pi|k|)^j k^2|\widehat f(k)|,
 \qquad
 8\pi^4t^2\sum_k(2\pi|k|)^j k^4|\widehat f(k)|.
\]
在 \(j=0\) 时导数因子按 \(1\) 取值。这些固定级数均收敛，且不依赖 \(t\)，从而得到两个 \(C^m\) 估计；第二式中的一次项正是 \(t(\mathcal Pf)''\)。
\end{solution}
